\section{Reference Engine, Controlled Tests, and Source-Linked Demonstrations}

\subsection{Validation ladder}

The release distinguishes six levels of evidence:
\begin{enumerate}
\item \textbf{L1---mechanical verification:} unit tests, schema validation, deterministic seeds, clean builds, manifests, and checksums;
\item \textbf{L2---controlled mechanism tests:} series generated by known linear, exponential, accelerating, saturating, and change-point processes;
\item \textbf{L3---source-linked historical demonstrations:} frozen excerpts from public primary-source datasets with explicit provenance and limitations;
\item \textbf{L4---preregistered prospective forecasts:} immutable cutoffs, inputs, intervals, triggers, and later scoring;
\item \textbf{L5---independent replication:} a separate team reproduces the software outputs and evaluates the method on unseen targets;
\item \textbf{L6---decision validation:} forecast-guided actions are evaluated against predeclared utility, loss, and regret.
\end{enumerate}
Version 2.0.0 supplies L1--L3 evidence and the registry, templates, and scoring code needed to accumulate L4--L6 evidence. It does not relabel controlled or retrospective demonstrations as prospective validation.

\subsection{End-to-end executable contract}

A single command now transforms a strict schema-valid input into a strict schema-valid probabilistic forecast, a self-contained executive HTML report, milestone crossing distributions, and reproducibility metadata:
\begin{Verbatim}[fontsize=\small,frame=single,breaklines=true]
fwta forecast protocol/v2/examples/canonical-reference-input.yaml \
  --schema protocol/v2/input.schema.json \
  --output results/ceiling/reference-forecast.json \
  --report results/ceiling/reference-forecast.html \
  --output-schema protocol/v2/output.schema.json
\end{Verbatim}
The output records the evidence clock, target semantics, six- and twelve-month pace, model fits and hindcasts, model weights, scenario assumptions and probabilities, interval-valued forecast points, milestones, bottlenecks, triggers, action plan, confidence, recalculation date, engine version, and input hash. A second, explicit 	exttt{fwta register} command can append a frozen output to the hash-chained registry. The release ships an empty live registry and a separately labelled demonstration chain, avoiding any implication that a bundled retrospective example was publicly preregistered.

\subsection{Controlled regime hindcasts}

The deterministic suite generates positive time series under known mechanisms and evaluates multiple candidate models across rolling origins. As expected, linear data favor the linear model and exponential data favor the exponential model. In accelerating and change-point settings, model averaging can perform best over the seeded horizon. In the logistic setting, an accelerating model can temporarily extrapolate well, illustrating a central warning: predictive success over a short window does not identify the true long-run mechanism.

\begin{table}[H]
\centering
\caption{Best rolling-origin model in each controlled generating regime, seed 20260728.}
\begin{tabular}{lll}
\toprule
Generating regime & Lowest-error model & Log RMSE \\
\midrule
accelerating & model average & 0.044 \\
change point & model average & 0.066 \\
exponential & exponential & 0.054 \\
linear & linear & 0.043 \\
logistic & accelerating & 0.058 \\
\bottomrule
\end{tabular}

\label{tab:synthetic-winners}
\end{table}

\begin{figure}[H]
\centering
\includegraphics[width=0.88\textwidth]{hindcast_heatmap.pdf}
\caption{Log-RMSE across controlled generating regimes and candidate models. Lower is better. The experiment tests pipeline behavior under known mechanisms, not real-world forecast accuracy.}
\label{fig:hindcast-heatmap}
\end{figure}

\subsection{Canonical-model ablations}

A controlled canonical system is generated with changing capability, cost, time compression, automation, parallelism, reliability, demand, adoption, and bottlenecks. Deliberately removing one or more layers increases forecast error in the seeded system. Because the components are correlated and generated according to the framework, these ablations are integrity tests rather than causal estimates or evidence that the same effect sizes hold in nature.

\begin{table}[H]
\centering
\caption{Controlled canonical-model ablations. Terminal ratio is predicted divided by observed at the final point.}
\begin{tabular}{lrr}
\toprule
Specification & Log RMSE & Terminal ratio \\
\midrule
full canonical & 0.035 & 0.94x \\
no reliability gain & 0.118 & 0.94x \\
no bottlenecks & 0.294 & 1.10x \\
no demand ceiling & 0.481 & 3.20x \\
no automation or parallelism & 0.685 & 0.92x \\
stale constant exponential & 1.334 & 18.98x \\
capability only & 1.459 & 0.19x \\
\bottomrule
\end{tabular}

\label{tab:ablation}
\end{table}

\begin{figure}[H]
\centering
\includegraphics[width=0.88\textwidth]{ablation_error.pdf}
\caption{Controlled ablation error. The exercise demonstrates expected failure modes when structural layers are omitted; it does not estimate real-world elasticities.}
\label{fig:ablation}
\end{figure}

\subsection{Source-linked historical demonstrations}

The first public demonstration uses 13 within-range observations transcribed from METR's Time Horizon 1.1 materials and excludes measurements beyond the stated 16-hour reliability boundary \citep{metr2026}. METR's metric is the human-expert duration of tasks that a system completes at a specified success probability; it is not agent runtime or economy-wide automation. The second demonstration uses seven selected frontier points from the official SWE-bench Verified leaderboard, whose benchmark contains human-validated software issues \citep{jimenez2024,swebench2026}. Both excerpts are frozen with source identifiers and caveats. They are intentionally small and should not be treated as complete mirrors, stable leaderboards, or prospective tests.

\begin{table}[H]
\centering
\caption{Source-linked demonstration summary. The interval shown is the base-scenario 80\% interval at the configured horizon.}
\begin{tabular}{p{0.28\textwidth}rp{0.12\textwidth}rrp{0.20\textwidth}}
\toprule
\textbf{Case} & \textbf{$n$} & \textbf{Best} & \textbf{Log-RMSE} & \textbf{Recent} & \textbf{Horizon median [P80]}\\
\midrule
METR Time Horizon 1.1 excerpt & 13 & change point & 0.292 & 1154.0\% & 960.0 [960.0, 960.0] \\
SWE-bench Verified frontier excerpt & 7 & logistic & 0.019 & 27.7\% & 83.5 [43.6, 86.1] \\
\bottomrule
\end{tabular}
\label{tab:empirical-cases}
\end{table}

\begin{figure}[H]
\centering
\includegraphics[width=0.94\textwidth]{empirical_cases.pdf}
\caption{Public-evidence demonstrations. METR extrapolation reaches a published measurement boundary; SWE-bench approaches a bounded percentage ceiling. The visualizations exercise different trajectory constraints and do not validate a general AI forecast.}
\label{fig:empirical-cases}
\end{figure}

The METR case illustrates why a measurement boundary must not be mistaken for a physical ceiling. The SWE-bench case illustrates why bounded outcomes require bounded models and why a leaderboard frontier mixes advances in base models, scaffolds, tools, budgets, and attempts. These examples test the data and reporting path; they do not establish causal attribution.

\subsection{Probabilistic rolling-origin scorecard}

Point error alone does not reveal whether uncertainty intervals are useful. We therefore apply six one-step rolling-origin forecasts to the compact within-range METR excerpt and score four predeclared model pools using log-RMSE, the continuous ranked probability score (CRPS), the weighted interval score (WIS), and empirical interval coverage. The sample is too small for general claims, but it exercises the probabilistic path under frozen historical cutoffs. In this run, the full model average has the lowest WIS and CRPS, while the accelerating-only pool has the lowest point RMSE. This disagreement is informative: the model with the best central prediction need not provide the best probabilistic forecast.

\begin{table}[H]
\centering
\caption{Probabilistic one-step rolling-origin hindcasts on the compact METR excerpt ($n=6$ forecast origins). Lower is better for RMSE, CRPS, and WIS. Coverage is descriptive at this sample size.}
\begin{tabular}{lrrrrr}
\toprule
Model pool & RMSE (log) & CRPS & WIS & 80\% cov. & 95\% cov. \\
\midrule
full model average & 0.298 & 45.8 & 138.1 & 83\% & 100\% \\
accelerating only & 0.285 & 49.1 & 145.7 & 100\% & 100\% \\
exponential only & 0.313 & 52.4 & 160.2 & 100\% & 100\% \\
change point only & 0.335 & 61.9 & 176.1 & 67\% & 100\% \\
\bottomrule
\end{tabular}

\label{tab:metr-probabilistic}
\end{table}

\begin{figure}[H]
\centering
\includegraphics[width=0.88\textwidth]{metr_probabilistic_hindcast.pdf}
\caption{Forecast medians, observed values, and 80\% intervals for the full model-average pool. The logarithmic vertical scale reflects the large dynamic range.}
\label{fig:metr-probabilistic}
\end{figure}

\subsection{Comparability audit: adoption is not one stable series}

The U.S. Census Bureau's Business Trends and Outlook Survey provides an instructive measurement-break case. Earlier releases measured AI use in producing goods or services, while the survey introduced questions about AI use in any business function beginning in November 2025 \citep{census2023ai,census2024btos,census2025history,census2026micro,census2026aiuse}. The broader construct produces a different level and cannot be concatenated mechanically with the earlier series. The release therefore stores the anchors in separate comparability groups and refuses to fit a single adoption trajectory without a validated bridge study.

\begin{table}[H]
\centering
\caption{BTOS AI-use anchors separated by measurement construct. The audit prevents a questionnaire change from being misclassified as adoption acceleration.}
\resizebox{\textwidth}{!}{\begin{tabular}{p{0.27\linewidth}p{0.28\linewidth}rrp{0.22\linewidth}}
\toprule
Comparability group & Construct & First & Last & Decision \\
\midrule
any\_business\_function\_v2 & AI used in any business function & 18.0\% & 19.8\% & Do not splice across the 2025 wording break. \\
goods\_services\_v1 & AI used to produce goods or services & 3.9\% & 10.0\% & Do not splice across the 2025 wording break. \\
\bottomrule
\end{tabular}
}
\label{tab:btos-break}
\end{table}

\subsection{Operational adoption and task compression evidence}

Technical possibility is not realized output, but operational indicators can reveal changing absorption and execution. GitHub reported a rise in average monthly merged pull requests from 35.0 million in 2024 to 43.2 million in 2025, more than one million pull requests created by its coding agent from May through September 2025, and a later change that reduced coding-agent start latency by a reported 50\% \citep{githuboctoverse2025,githubagent2026}. These figures have different denominators and cannot be multiplied into a single productivity index. They are included as evidence that adoption volume and workflow latency have their own clocks and require separate measurement.

\begin{table}[H]
\centering
\caption{Selected operational evidence records retained with source and unit metadata.}
\begin{tabular}{p{0.20\textwidth}p{0.38\textwidth}rp{0.18\textwidth}}
\toprule
\textbf{Period} & \textbf{Metric} & \textbf{Value} & \textbf{Unit}\\
\midrule
2024 & merged\_pull\_requests\_monthly\_average & 35,000,000 & pull requests/month \\
2025 & merged\_pull\_requests\_monthly\_average & 43,200,000 & pull requests/month \\
2025-05 to 2025-09 & coding-agent pull requests created & 1,000,000 & pull requests \\
2026-03-19 & coding-agent start-time improvement & 50\% & fraction faster \\
\bottomrule
\end{tabular}
\label{tab:github-evidence}
\end{table}

\subsection{What the release proves---and does not prove}

The release demonstrates that the equations, schemas, engine, report, registry, controlled experiments, and public-data adapters can execute coherently and reproducibly in the documented environment. It demonstrates that different modelling choices produce materially different forecasts under known mechanisms and on small source-linked series. It does not prove that the canonical model is universally superior, that scenario probabilities are calibrated, that a discontinuity will occur, or that a particular future forecast is correct. Those claims require the prospective and independent program in Section~\ref{sec:validation-program}.
